\subsection{若干引理}

引理 1．——设 A 是诺特半局部环，B 是有限 A-代数。则环 B 是半局部且诺特的；设 \(\mathfrak{m}_1, \ldots, \mathfrak{m}_n\) 是其极大理想。

从 B 到 \(\prod_{i=1}^{n} \hat{\mathbf{B}}_{\mathfrak{m}_i}\) 的典范同态可延拓为从 \(\hat{\mathbf{A}} \otimes_{\mathbf{A}} \mathbf{B}\) 到 \(\prod_{i=1}^{n} \hat{\mathbf{B}}_{\mathfrak{m}_i}\) 的同构。

根据 IV，§ 2，第 5 号，命题 9 的推论 3，环 B 是半局部的，且 \(\mathfrak{m}_{A}B\) 是它的定义理想。根据 III，§ 3，第 4 号，定理 3 (ii)，环 \(\hat{\mathbf{A}} \otimes_{\mathbf{A}} \mathbf{B}\) 是 B 关于由其根基定义的拓扑的完备化；然后应用 III，§ 2，第 13 号，命题 19 的推论。

引理 2．——设 A 是诺特环，M 是一个 A-模。典范映射从 M 到乘积 \(\prod_{\mathfrak{p} \in \operatorname{Ass}(M)} M_{\mathfrak{p}}\) 是单射。

事实上，取 M 的非零元素 \(m\)；则 Ann( \(m\) ) 包含于 A 的某个与 M 相关的素理想 \(\mathfrak{p}\)（IV，§ 1，第 1 号，命题 2），且 \(m\) 在 \(M_{\mathfrak{p}}\) 中的像非零（II，§ 2，第 2 号，命题 4）。

引理 3．——设 A 是诺特环，x 是 A 的元素，M 是有限生成的 A-模，且 \(\mathfrak{p}\) 是 M 的一个相关素理想。设倍乘映射 \(x_{\mathbf{M}}\) 是单射。令 \(\mathfrak{q}\) 是 A 的一个素理想，在包含 \(\mathfrak{p} + xA\) 的素理想中极小。则 \(\mathfrak{q}\) 是 A-模 M/xM 的相关素理想。

令 N 表示 M 的子模，由满足 \(\mathfrak{p}m = 0\) 的元素 m 组成。我们有 N ∩ xM = xN；事实上，若 M 的一个元素 m 满足 \(\mathfrak{p}xm = 0\)，则有 \(\mathfrak{p}m = 0\)，因为 \(x_{M}\) 是单射，从而 \(m \in N\)。因此，A-模 N/xN 同构于 M/xM 的子模 (N + xM)/xM，所以只需证明 \(\mathfrak{q}\) 是 N/xN 的相关素理想。由于 \(\mathfrak{p}\) 与 M 相关，存在 M 的元素 m，使 \(\mathfrak{p} = \operatorname{Ann}(m)\)；有 \(m \in N\)，故 \(\mathfrak{p} = \operatorname{Ann}(N)\)，于是 \(\operatorname{Supp}(N/xN) = V(\mathfrak{p} + xA)\)（II，§ 4，第 4 号，命题 18 的推论）；因此，\(\mathfrak{q}\) 是 N/xN 的相关素理想（IV，§ 1，第 4 号，定理 2）。

引理 4．——设 A 是离散赋值环，B 是诺特局部环，且 \(\rho: \mathbf{A} \to \mathbf{B}\) 是局部且平坦的同态。若环 \(\kappa_{\mathbf{A}} \otimes_{\mathbf{A}} \mathbf{B}\) 是约化的，则 B 是约化的。

假设 B 中存在非零幂零元素 \(x\)，并令 \(\pi\) 为 A 的赋值参数。由于 \(\pi \mathbf{B} \subset \mathfrak{m}_{\mathbf{B}}\)，环 B 关于 \(\pi \mathbf{B}\) -进拓扑是分离的。因此存在 \(n \in \mathbf{N}\) 和 \(y \in \mathbf{B}\)，使 \(x = \pi^n y\) 且 \(y \notin \pi \mathbf{B}\)。由于 B 平坦于 A，乘以 \(\pi\) 的映射在 B 中是单射。因此，\(y\) 在 \(\mathbf{B}/\pi \mathbf{B}\) 中的类是非零幂零元素，这与假设矛盾。
% semantic-fragments: legacy-input-seam
\relax{}
